% ============ 第八章 §3 赋值的扩充（完备情形） (原书 pp.287-294) ============
\section{赋值的扩充（完备情形）}\label{sec:ch8-s3}

设 $E/F$ 是域的有限次扩张, 我们要讨论的问题是: $F$ 的每个素除子 $P$ 到 $E$ 有哪些扩充素除子？本节研究 $(F,P)$ 是完备赋值域的情形. 对于这种情形, 结论很简单: $P$ 到 $E$ 只有惟一的扩充 $P_E$, $(E,P_E)$ 仍是完备赋值域, 并且 $P_E$ 可以有明确的刻画.

如果 $P$ 是阿基米德素除子, 由定理 \ref{thm:8.3} 知 $(F,P)$ 为 $(\R,\infty)$ 或者 $(\C,\infty)$. 由于 $\C$ 是代数封闭的, 所以 $E/F$ 只能是 $\R/\R$, $\C/\R$ 或者 $\C/\C$. 并且 $E$ 和 $F$ 的素除子均为 $\infty$, 所以上述结论成立. 以下设 $(F,P)$ 是完备的非阿赋值域.

我们可以像数域扩张的情形那样, 对于域的任意有限扩张 $E/F$ 定义范映射 $N=N_{E/F}$. 设 $[E:F]=n$. 取向量空间 $E$ 的一组 $F$-基 $\omega_1,\cdots,\omega_n$, 则 $E$ 中每个元素都惟一表示成 $a_1\omega_1+\cdots+a_n\omega_n$（$a_i\in F$）. 于是对每个 $\alpha\in E$,
\[
  \alpha\left(\begin{array}{c}\omega_1\\ \vdots\\ \omega_n\end{array}\right)
  =M\left(\begin{array}{c}\omega_1\\ \vdots\\ \omega_n\end{array}\right),
\]
其中 $M$ 是元素属于 $F$ 的 $n$ 阶方阵. 定义 $N(\alpha)=N_{E/F}(\alpha)$ 为 $M$ 的行列式 $|M|$. 易知这个定义与基 $\omega_1,\cdots,\omega_n$ 的选取方式无关. 并且范映射 $N=N_{E/F}\colon E\to F$ 有如下性质:

(1) $N(\alpha)=0\Longleftrightarrow\alpha=0$;

(2) $N(\alpha\beta)=N(\alpha)N(\beta)$;

(3) 若 $\alpha\in F$, 则 $N(\alpha)=\alpha^n$.

如果 $F$（从而 $E$）是特征零域, 则 $E/F$ 是单扩张, 即 $E=F(\theta)$. 令 $f(x)\in F[x]$ 是 $\theta$ 在 $F$ 上的最小多项式, 则 $f(x)$ 为 $F[x]$ 中 $n$ 次不可约多项式. $f(x)$ 在 $F$ 的代数闭包中有 $n$ 个不同的根 $\theta_1,\cdots,\theta_n$. 这时可以证明范映射可以表成另一种形式: 由于 $1,\theta,\cdots,\theta^{n-1}$ 是 $E$ 对 $F$ 的一组基, 从而 $E$ 中元素可惟一表成 $\alpha=a_0+a_1\theta+\cdots+a_{n-1}\theta^{n-1}$. 我们有
\[
  N_{E/F}(\alpha)=\prod_{i=1}^{n}\left(a_0+a_1\theta_i+\cdots+a_{n-1}\theta_i^{n-1}\right),
\]
这个表达式和代数数域的情形一致.

\begin{theorem}\label{thm:8.13}
（存在性）设 $(F,P)$ 为完备非阿赋值域, $O_F$ 为 $F$ 对 $P$ 的赋值环, $\varphi\in P$, $E/F$ 为域的 $n$ 次扩张, $N=N_{E/F}$.

(1) 对每个 $\alpha\in E$, 令
\[
  \varphi_E(\alpha)=\varphi(N(\alpha))^{\frac{1}{n}}\geqslant 0,
\]
则 $\varphi_E$ 是 $E$ 的（非阿）赋值, 并且是 $\varphi$ 的扩充（即 $\varphi_E|_F=\varphi$）.

(2) $\varphi$ 为离散赋值 $\Leftrightarrow\varphi_E$ 为离散赋值.

(3) 设 $\varphi_E$ 的素除子为 $P_E$, $E$ 对 $P_E$ 的赋值环为 $O_E$, 则 $O_E$ 是 $O_F$ 在 $E$ 中的整闭包, 并且 $O_F=O_E\cap F$.
\end{theorem}

\begin{proof}
(1) 当 $\alpha\in F$ 时, $\varphi_E(\alpha)=\varphi(N(\alpha))^{\frac{1}{n}}=\varphi(\alpha^n)^{\frac{1}{n}}=\varphi(\alpha)$. 所以 $\varphi_E|_F=\varphi$. 再证 $\varphi_E$ 为 $E$ 的赋值. 易知
\[
  \varphi_E(\alpha)=0\Leftrightarrow N(\alpha)=0\Leftrightarrow\alpha=0,
\]
\[
  \varphi_E(\alpha\beta)=\varphi_E(\alpha)\varphi_E(\beta),
\]
所以只需证明对每个 $\alpha\in E$, $\varphi_E(\alpha)\leqslant 1\Rightarrow\varphi_E(1+\alpha)\leqslant 1$.

设 $f(x)=x^d+a_1x^{d-1}+\cdots+a_d\in F[x]$ 是 $\alpha$ 在 $F$ 上的最小多项式, 则 $N(\alpha)=\pm a_d^m$, $m=\frac{n}{d}=[E:F(\alpha)]$. 于是
\[
  \varphi_E(\alpha)\leqslant 1\Leftrightarrow\varphi(N(\alpha))\leqslant 1\Leftrightarrow\varphi(a_d)\leqslant 1\Leftrightarrow a_d\in O_F.
\]
由于 $f(x)$ 在 $F[x]$ 中不可约, 由定理 \ref{thm:8.7} 可知 $a_1,\cdots,a_d$ 均属于 $O_F$, 即 $f(x)\in O_F[x]$. 注意 $g(x)=f(x-1)\in O_F[x]$ 是 $\alpha+1$ 在 $F$ 上的最小多项式, 而
\[
  N(1+\alpha)=(\pm g(0))^m=(\pm f(1))^m\in O_F.
\]
从而 $\varphi_E(1+\alpha)\leqslant 1$. 这就表明 $\varphi_E$ 是非阿赋值.

(2) 由于 $\log\varphi_E(\alpha)=\frac{1}{n}\log\varphi(N(\alpha))$, 可知此式一边取离散值, 则另一边亦然.

(3) 我们在证明 (1) 时已得到:
\[
  \alpha\in O_E\Leftrightarrow\varphi_E(\alpha)\leqslant 1\Leftrightarrow\alpha\ \text{在 } F\ \text{上的最小多项式属于 } O_F[x]
  \Leftrightarrow\alpha\ \text{在 } O_F\ \text{上整}.
\]
从而 $O_E$ 为 $O_F$ 在 $E$ 中的整闭包. 最后 $O_F=F\cap O_E$ 显然成立.
\end{proof}

现在证明素除子扩充的惟一性. 我们知道, $E$ 的不同素除子给出 $E$ 的不同拓扑结构. 为证扩充的惟一性, 只需证明 $F$ 的 $\varphi$-拓扑到 $E$ 的所有扩充都给出 $E$ 的同一拓扑即可.

\begin{definition}\label{def:ch8-2}
设 $(F,\varphi)$ 是非阿赋值域. $F$ 上的向量空间 $X$ 叫作对 $\varphi$ 的\emph{赋范空间}, 是指存在映射 $\|\cdot\|\colon X\to\R$ 满足（对 $a\in F$, $\xi,\eta\in X$）:

(i) $\|\xi\|\geqslant 0$, 并且 $\|\xi\|=0\Leftrightarrow\xi=0$;

(ii) $\|a\xi\|=\varphi(a)\|\xi\|$;

(iii) $\|\xi+\eta\|\leqslant\|\xi\|+\|\eta\|$.

这时, $\|\xi\|$ 叫作 $\xi$ 的范.
\end{definition}

\begin{example}
设 $(E,\varphi_E)\supset(F,\varphi)$ 是非阿赋值域的扩张. 令 $\varphi_E=\|\cdot\|$, 则 $E$ 是 $F$ 上对 $\varphi$ 的赋范空间.
\end{example}

\begin{example}
设 $(F,\varphi)$ 是非阿赋值域, 则 $F$ 上每个有限维向量空间 $X$ 均可如下作成 $(F,\varphi)$ 上的赋范空间. 取 $X$ 的一组 $F$-基 $\omega_1,\cdots,\omega_n$. 对 $X$ 中元素 $\alpha=\sum_{i=1}^{n}a_i\omega_i$（$a_i\in F$）定义
\[
  \|\xi\|=\max_{1\leqslant i\leqslant n}\{\varphi(a_i)\}.
\]
不难验证 $\|\cdot\|$ 是 $X$ 的范. 并且若 $(F,\varphi)$ 完备, 则 $X$ 对于由 $\|\cdot\|$ 所诱导的拓扑是完备拓扑空间.
\end{example}

下面是泛函分析中的一个熟知结果.

\begin{lemma}\label{lem:ch8-2}
设 $(F,\varphi)$ 是完备非阿赋值域, $X$ 是 $F$ 上有限维（$n$ 维）向量空间. $\|\cdot\|$ 是例 $2$ 中定义的范, $|\cdot|$ 是 $X$ 对于 $(F,\varphi)$ 的任意范, 则存在常数 $D_1,D_2>0$, 使得对每个 $\xi\in X$,
\[
  D_1\|\xi\|\leqslant|\xi|\leqslant D_2\|\xi\|,
\]
并且 $X$ 对于 $|\cdot|$ 也是完备的.

（换句话说, $X$ 对于 $(F,\varphi)$ 的任意两个范诱导出 $X$ 上同样的拓扑.）
\end{lemma}

\begin{proof}
取 $D_2=|\omega_1|+\cdots+|\omega_n|$, 则由三角形不等式知
\[
  |\xi|\leqslant\sum_{i=1}^{n}\varphi(a_i)|\omega_i|\leqslant\max_{1\leqslant i\leqslant n}(\varphi(a_i))D_2=D_2\|\xi\|.
\]
下面对 $n$ 归纳证明 $D_1>0$ 的存在性. 当 $n=1$ 时, $X=F$, 命题显然成立. 现设命题对 $n-1$ 成立. 以 $Y_i$ 表示 $\omega_1,\cdots,\omega_{i-1},\omega_{i+1},\cdots,\omega_n$ 张成的 $X$ 的子空间, 则 $|\cdot|$ 也为 $Y_i$ 上的范. 由归纳假设, $Y_i$ 对于 $|\cdot|$ 完备, 从而 $Y_i$ 为 $X$ 的闭子空间, 于是 $\bigcup_{i=1}^{n}(Y_i+\omega_i)$ 是 $X$ 的闭集. 但是 $0\notin Y_i+\omega_i\ (1\leqslant i\leqslant n)$, 所以存在 $D_1>0$, 使得
\[
  |\eta_i+\omega_i|\geqslant D_1\qquad(\text{对每个 } \eta_i\in Y_i,\ 1\leqslant i\leqslant n).
\]
设 $0\neq\alpha=a_1\omega_1+\cdots+a_n\omega_n\in X$. 令 $\varphi(a_r)=\max\limits_{1\leqslant i\leqslant n}\{\varphi(a_i)\}$, 则 $0\neq\varphi(a_r)=\|\alpha\|$, $a_r\neq 0$, 于是 $a_r^{-1}\alpha\in Y_r+\omega_r$, 从而
\[
  |\alpha a_r^{-1}|\geqslant D_1,\qquad|\alpha|\geqslant\varphi(a_r)D_1=D_1\|\alpha\|,
\]
因此命题对 $n$ 成立. 这表明 $\|\cdot\|$ 和 $|\cdot|$ 是等价的拓扑. 由 $X$ 对 $\|\cdot\|$ 完备可知对 $|\cdot|$ 也完备.
\end{proof}

\begin{theorem}\label{thm:8.14}
（惟一性）设 $(F,\varphi)$ 是完备非阿赋值域, 则定理 \ref{thm:8.13} 中的 $\varphi_E$ 是 $\varphi$ 到 $E$ 的惟一扩充, 并且 $(E,\varphi_E)$ 是完备的.
\end{theorem}

\begin{proof}
设 $\varphi_1$ 和 $\varphi_2$ 是 $\varphi$ 到 $E$ 的两个扩充, 则它们都是向量空间 $E$ 对于 $(F,\varphi)$ 的范. 由引理 \ref{lem:ch8-2} 可知 $E$ 的 $\varphi_1$-拓扑和 $\varphi_2$-拓扑是等价的, 从而 $\varphi_1$ 和 $\varphi_2$ 等价. 于是存在 $s>0$, 使得 $\varphi_1=\varphi_2^s$. 但是
\[
  \varphi=\varphi_1|_F=\varphi_2^s|_F=\varphi_2|_F^s=\varphi^s,
\]
可知 $s=1$, 即 $\varphi_1=\varphi_2$. 最后由引理 \ref{lem:ch8-2} 知 $E$ 对于 $\varphi_E$ 拓扑是完备的.
\end{proof}

\begin{theorem}\label{thm:8.15}
设 $(F,\varphi)$ 是完备赋值域, $E/F$ 为有限伽罗瓦扩张, $\varphi_E$ 是 $\varphi$ 到 $E$ 的（惟一）扩充. 如果 $E$ 中元素 $\alpha$ 和 $\beta$ 是 $F$-共轭的, 则 $\varphi_E(\alpha)=\varphi_E(\beta)$.
\end{theorem}

\begin{proof}
$\alpha$ 和 $\beta$ 是 $F$-共轭元素, 是指存在 $E$ 的 $F$-自同构 $\sigma$, 使得 $\sigma(\alpha)=\beta$. 由于 $\varphi_E\circ\sigma$ 也是 $E$ 的赋值并且为 $\varphi$ 的扩充. 由惟一性知 $\varphi_E\circ\sigma=\varphi_E$, 即
\[
  \varphi_E(\alpha)=\varphi_E(\sigma(\alpha))=\varphi_E(\beta).
\]
\end{proof}

\begin{remark}
设 $(F,\varphi)$ 是完备赋值域, $\Omega$ 是 $F$ 的代数闭包, 则 $\varphi$ 到 $\Omega$ 存在惟一扩充. 这是因为: 对每个 $\alpha\in\Omega$, 均存在 $F$ 的有限扩域 $E$, 使得 $\alpha\in E$（如取 $E=F(\alpha)$）. 令 $\varphi_E$ 是 $\varphi$ 到 $E$ 的惟一扩充, 定义 $\varphi_\Omega(\alpha)=\varphi_E(\alpha)$. 不难证明 $\varphi_\Omega(\alpha)$ 与域 $E$ 的取法无关. $\varphi_\Omega$ 为域 $\Omega$ 的赋值, 并且是 $\varphi$ 到 $\Omega$ 的惟一扩充.
\end{remark}

下面是完备离散赋值域扩张的基本结果.

\begin{theorem}\label{thm:8.16}
设 $(F,P)\subseteq(E,Q)$ 是完备离散赋值域的有限扩张, $n=[E:F]$. 令 $v_P$ 和 $v_Q$ 分别是素除子 $P$ 和 $Q$ 中的标准指数赋值（即 $v_P(F^{\times})=v_Q(E^{\times})=\Z$）, 则

(1) 对于 $a\in F$, $v_Q(a)=e(Q/P)v_P(a)$;

(2) $e(Q/P)f(Q/P)=n$;

(3) 对于 $\alpha\in E$, $v_P(N_{E/F}(\alpha))=f(Q/P)v_Q(\alpha)$.
\end{theorem}

\begin{proof}
简记 $e=e(Q/P)$, $f=f(Q/P)$, $N=N_{E/F}$.

(1) 令 $v'=v_Q|_F$, 则 $v'$ 和 $v_P$ 等价, 所以 $v'=sv_P\ (s>0)$. 设 $\pi$ 为 $F$ 中的素元, 则 $v'(\pi)=sv_P(\pi)=s\in v'(F^{\times})$, 并且 $v'(F^{\times})=s\Z$. 但是
\[
  e=[v_Q(E^{\times}):v_Q(F^{\times})]=[\Z:v'(F^{\times})]=[\Z:s\Z]=s,
\]
因此 $v'=ev_P$.

(2) 以 $O_E$ 和 $O_F$ 分别表示 $E$ 和 $F$ 的赋值环, $\overline{E}$ 和 $\overline{F}$ 为相应的剩余类域. 令 $\omega_1,\cdots,\omega_f\in O_E$, 使得 $\bar{\omega}_1,\cdots,\bar{\omega}_f$ 是 $\overline{E}$ 的一组 $\overline{F}$-基, 令 $\pi$ 为 $E$ 中的素元, 则 $1,\pi,\pi^2,\cdots,\pi^{e-1}$ 为 $v_Q(E^{\times})$ 对于 $v_Q(F^{\times})$ 的陪集完全代表系. 我们证明
\begin{equation}\tag{*}
  O_E=\bigoplus_{\substack{1\leqslant i\leqslant f\\ 0\leqslant j\leqslant e-1}}\omega_i\pi^jO_F\qquad(\text{加法群的直和}).
\end{equation}
由于 $E$ 和 $F$ 分别是 $O_E$ 和 $O_F$ 的商域, 可由此推出 $E=\bigoplus_{i,j}\omega_i\pi^jF$. 于是 $ef=[E:F]=n$.

现在证明直和分解 $(*)$. 令 $L=\sum_{i=1}^{f}\omega_iO_F$, 则 $O_E=L+\pi O_E$（$\supseteq$: 显然; $\subseteq$: 令 $\alpha\in O_E$, 则 $\bar{\alpha}=\bar{a}_1\bar{\omega}_1+\cdots+\bar{a}_f\bar{\omega}_f$, $a_i\in O_F$, 于是 $\alpha-(a_1\omega_1+\cdots+a_f\omega_f)\in\pi O_E$）. 从而
\[
  O_E=L+\pi(L+\pi O_E)=L+\pi L+\pi^2O_E
  =L+\pi L+\cdots+\pi^{e-1}L+\pi^eO_E.
\]
令 $M=L+\pi L+\cdots+\pi^{e-1}L=\sum_{i,j}\omega_i\pi^jO_F$. 由于 $\omega_i\pi^jO_F$ 中元素 $\alpha$ 均有 $v_Q(\alpha)\equiv j\ (\mathrm{mod}\ e)$, 所以 $M=\bigoplus_{i,j}\omega_i\pi^jO_F$ 为直和. 我们只需再证 $M=O_E$. 设 $\pi_F$ 是 $F$ 中素元, 则
\begin{align*}
  O_E&=M+\pi^eO_E=M+\pi_FO_E\\
  &=M+\pi_FM+\cdots+\pi_F^{r-1}M+\pi_F^rO_E\qquad(r\geqslant 1).\tag{**}
\end{align*}
但是对每个 $l\geqslant 1$, $\pi_F^lL\subseteq L$, 从而
\[
  \pi_F^lM=\pi_F^l(L+\pi L+\cdots+\pi^{e-1}L)\subseteq M.
\]
再由 $(**)$ 式可知 $O_E=M+\pi_F^rO_E$（对每个 $r\geqslant 1$）. 这表明 $M$ 在 $O_E$ 中稠密. 另一方面, $O_F$ 为 $F$ 中闭集, 可知 $M=\bigoplus_{i,j}\omega_i\pi^jO_F$ 为 $O_E$ 中闭集. 于是 $M=O_E$.

(3) 设 $v_Q=-\log\circ\varphi_Q$, 其中 $\varphi_Q$ 为 $Q$ 中赋值. 令 $\varphi_Q|_F=\varphi_P$, 则在 $F$ 上有 $ev_P=-\log\circ\varphi_P$. 但是 $\varphi_Q(\alpha)=\varphi_P(N(\alpha))^{\frac{1}{n}}$. 因此
\begin{align*}
  v_Q(\alpha)&=\frac{1}{n}\left(-\log\circ\varphi_P(N(\alpha))\right)\\
  &=\frac{e}{n}v_P(N(\alpha))=\frac{1}{f}v_P(N(\alpha)).
\end{align*}
\end{proof}

\begin{Exercises}
\item 设 $E$ 是 $\Q_p$ 的有限扩张, $v_p$ 为 $\Q_p$ 的标准指数赋值, $v$ 为 $v_p$ 到 $E$ 的惟一扩充. 求证:

(1) 对于 $x\in E$, 当 $v(x)>\frac{1}{p-1}$ 时, 级数
\[
  \exp(x)=1+x+\frac{x^2}{2!}+\cdots+\frac{x^n}{n!}+\cdots
\]
收敛, 并且当 $n\geqslant 2$ 时 $v\left(\frac{x^n}{n!}\right)>v(x)$. 于是
\[
  v(x)=v(\exp(x)-1).
\]
(2) 若 $x,y\in E$, $v(x)>\frac{1}{p-1}$, $v(y)>\frac{1}{p-1}$, 则 $\exp(x+y)=\exp(x)\exp(y)$.

(3) 当 $x\in E$, $v(x)>0$ 时, 级数
\[
  \log(1+x)=x-\frac{x^2}{2}+\frac{x^3}{3}-\cdots+(-1)^{n+1}\frac{x^n}{n}+\cdots
\]
收敛. 如果 $v(x)>\frac{1}{p-1}$, 则当 $n\geqslant 2$ 时, $v\left(\frac{x^n}{n}\right)>v(x)$. 于是 $v(\log(1+x))=v(x)$.

(4) 若 $x,y\in E$, $v(x-1)>0$, $v(y-1)>0$, 则
\[
  \log(xy)=\log x+\log y.
\]
(5) 若 $x\in E$, $v(x)>\frac{1}{p-1}$, 则
\[
  \exp(\log(1+x))=1+x,\qquad\log(\exp(x))=x.
\]
特别当 $r\geqslant\left[\frac{e}{p-1}\right]+1$ 时, 映射
\[
  \log\colon U_r\to\mathfrak{p}_E^r,\qquad\exp\colon\mathfrak{p}_E^r\to U_r=1+\mathfrak{p}_E^r
\]
给出乘法群 $U_r$ 和加法群 $\mathfrak{p}_E^r$ 之间互逆的群同构, 其中 $\mathfrak{p}_E$ 是 $E$ 的素理想, 而 $e=e(\mathfrak{p}_E/p)$.

(6) 群同态 $f\colon U_1=1+\mathfrak{p}_E\to E$, $x\mapsto\log x$ 的核为
\[
  \operatorname{Ker}f=\{\alpha\in E\mid\text{存在 } n\geqslant 0,\ \text{使得 } \alpha^{p^n}=1\}.
\]

\item 对于 $\Q_p$ 的所有二次扩域 $E$, 计算 $E/\Q_p$ 的分歧指数和剩余类域次数.

\item 设 $(F,P)\subseteq(E,Q)$ 是完备离散赋值域的有限扩张. 求证:

(1) $\mathfrak{p}_FO_E=\mathfrak{p}_E^e$, $e=e(Q/P)$;

(2) 对于整数 $m\geqslant 0$, 试问 $O_F$ 中理想 $\mathfrak{p}_E^m\cap O_F=?$

(3) 若 $A$ 为 $O_E$ 中的理想, 如何刻画由集合 $N(A)=\{N_{E/F}(\alpha)\mid\alpha\in A\}$ 在 $O_F$ 中生成的理想？
\end{Exercises}
